% Môn: Toán
% Lớp: 10
% Chương: Chương II. Bất phương trình và hệ bất phương trình bậc nhất hai ẩn
% Bài: Bài 3. Bất phương trình bậc nhất hai ẩn
% Loại: Dạng mẫu
% File riêng, không trộn vào de.tex
% Định nghĩa lệnh Lý thuyết — dùng khi biên dịch lt.tex bằng TeX.
% App web đọc lt.tex trực tiếp (bỏ \input file này).
% Trong lt.tex, từ thư mục bài:
%   % Định nghĩa lệnh Lý thuyết — dùng khi biên dịch lt.tex bằng TeX.
% App web đọc lt.tex trực tiếp (bỏ \input file này).
% Trong lt.tex, từ thư mục bài:
%   % Định nghĩa lệnh Lý thuyết — dùng khi biên dịch lt.tex bằng TeX.
% App web đọc lt.tex trực tiếp (bỏ \input file này).
% Trong lt.tex, từ thư mục bài:
%   \input{../../../../_lenh/lythuyet.tex}
% từ gốc repo:
%   \input{ngan-hang/_lenh/lythuyet.tex}
\makeatletter
\providecommand{\indam}[1]{\textbf{#1}}
\providecommand{\chuy}[1]{\par\smallskip\noindent\textit{#1}\par}
\providecommand{\luuy}[1]{\par\smallskip\noindent\textbf{Lưu ý.} #1\par}
\providecommand{\ghichu}[1]{\par\smallskip\noindent\textbf{Ghi chú.} #1\par}
\providecommand{\vidu}[1]{\par\smallskip\noindent\textbf{Ví dụ.} #1\par}
\providecommand{\Blythuyet}{\subsection*{Lý thuyết}}
% Hai cột: kiến thức | hình
\providecommand{\haicotchay}[2]{%
  \par\medskip\noindent
  \begin{minipage}[t]{0.52\linewidth}#1\end{minipage}\hfill
  \begin{minipage}[t]{0.46\linewidth}#2\end{minipage}%
  \par\medskip
}
\providecommand{\kienthuccot}[2]{\haicotchay{#1}{#2}}
% Dự phòng nếu chưa nạp graphicx / caption (không ghi đè bản thật)
\providecommand{\resizebox}[3]{#3}
\providecommand{\captionof}[2]{\par\smallskip\noindent\textit{#2}\par}
\newenvironment{hoatdong}[1]{%
  \par\medskip\noindent\textbf{Hoạt động.} #1\par\smallskip
}{\par\medskip}
\newenvironment{emcobiet}{%
  \par\medskip\noindent\textbf{Em có biết}\par\smallskip
}{\par\medskip}
\newenvironment{emdahoc}{%
  \par\medskip\noindent\textbf{Em đã học}\par\smallskip
}{\par\medskip}
\newenvironment{emcothe}{%
  \par\medskip\noindent\textbf{Em có thể}\par\smallskip
}{\par\medskip}
\newenvironment{kienthuc}{%
  \par\medskip\noindent\textbf{Kiến thức cốt lõi}\par\smallskip
}{\par\medskip}
\newenvironment{traloi}{%
  \par\smallskip\noindent\textit{Trả lời / ghi chú của em}\par
  \vspace{1.2em}%
}{\par\medskip}
\newenvironment{phuongphap}{%
  \par\medskip\noindent\textbf{Phương pháp giải}\par\smallskip
}{\par\medskip}
\newenvironment{vidumau}{%
  \par\medskip\noindent\textbf{Bài mẫu}\par\smallskip
}{\par\medskip}
\newenvironment{dangmau}[1]{%
  \par\medskip\noindent\textbf{Dạng mẫu.} #1\par\smallskip
}{\par\medskip}
\makeatother

% từ gốc repo:
%   % Định nghĩa lệnh Lý thuyết — dùng khi biên dịch lt.tex bằng TeX.
% App web đọc lt.tex trực tiếp (bỏ \input file này).
% Trong lt.tex, từ thư mục bài:
%   \input{../../../../_lenh/lythuyet.tex}
% từ gốc repo:
%   \input{ngan-hang/_lenh/lythuyet.tex}
\makeatletter
\providecommand{\indam}[1]{\textbf{#1}}
\providecommand{\chuy}[1]{\par\smallskip\noindent\textit{#1}\par}
\providecommand{\luuy}[1]{\par\smallskip\noindent\textbf{Lưu ý.} #1\par}
\providecommand{\ghichu}[1]{\par\smallskip\noindent\textbf{Ghi chú.} #1\par}
\providecommand{\vidu}[1]{\par\smallskip\noindent\textbf{Ví dụ.} #1\par}
\providecommand{\Blythuyet}{\subsection*{Lý thuyết}}
% Hai cột: kiến thức | hình
\providecommand{\haicotchay}[2]{%
  \par\medskip\noindent
  \begin{minipage}[t]{0.52\linewidth}#1\end{minipage}\hfill
  \begin{minipage}[t]{0.46\linewidth}#2\end{minipage}%
  \par\medskip
}
\providecommand{\kienthuccot}[2]{\haicotchay{#1}{#2}}
% Dự phòng nếu chưa nạp graphicx / caption (không ghi đè bản thật)
\providecommand{\resizebox}[3]{#3}
\providecommand{\captionof}[2]{\par\smallskip\noindent\textit{#2}\par}
\newenvironment{hoatdong}[1]{%
  \par\medskip\noindent\textbf{Hoạt động.} #1\par\smallskip
}{\par\medskip}
\newenvironment{emcobiet}{%
  \par\medskip\noindent\textbf{Em có biết}\par\smallskip
}{\par\medskip}
\newenvironment{emdahoc}{%
  \par\medskip\noindent\textbf{Em đã học}\par\smallskip
}{\par\medskip}
\newenvironment{emcothe}{%
  \par\medskip\noindent\textbf{Em có thể}\par\smallskip
}{\par\medskip}
\newenvironment{kienthuc}{%
  \par\medskip\noindent\textbf{Kiến thức cốt lõi}\par\smallskip
}{\par\medskip}
\newenvironment{traloi}{%
  \par\smallskip\noindent\textit{Trả lời / ghi chú của em}\par
  \vspace{1.2em}%
}{\par\medskip}
\newenvironment{phuongphap}{%
  \par\medskip\noindent\textbf{Phương pháp giải}\par\smallskip
}{\par\medskip}
\newenvironment{vidumau}{%
  \par\medskip\noindent\textbf{Bài mẫu}\par\smallskip
}{\par\medskip}
\newenvironment{dangmau}[1]{%
  \par\medskip\noindent\textbf{Dạng mẫu.} #1\par\smallskip
}{\par\medskip}
\makeatother

\makeatletter
\providecommand{\indam}[1]{\textbf{#1}}
\providecommand{\chuy}[1]{\par\smallskip\noindent\textit{#1}\par}
\providecommand{\luuy}[1]{\par\smallskip\noindent\textbf{Lưu ý.} #1\par}
\providecommand{\ghichu}[1]{\par\smallskip\noindent\textbf{Ghi chú.} #1\par}
\providecommand{\vidu}[1]{\par\smallskip\noindent\textbf{Ví dụ.} #1\par}
\providecommand{\Blythuyet}{\subsection*{Lý thuyết}}
% Hai cột: kiến thức | hình
\providecommand{\haicotchay}[2]{%
  \par\medskip\noindent
  \begin{minipage}[t]{0.52\linewidth}#1\end{minipage}\hfill
  \begin{minipage}[t]{0.46\linewidth}#2\end{minipage}%
  \par\medskip
}
\providecommand{\kienthuccot}[2]{\haicotchay{#1}{#2}}
% Dự phòng nếu chưa nạp graphicx / caption (không ghi đè bản thật)
\providecommand{\resizebox}[3]{#3}
\providecommand{\captionof}[2]{\par\smallskip\noindent\textit{#2}\par}
\newenvironment{hoatdong}[1]{%
  \par\medskip\noindent\textbf{Hoạt động.} #1\par\smallskip
}{\par\medskip}
\newenvironment{emcobiet}{%
  \par\medskip\noindent\textbf{Em có biết}\par\smallskip
}{\par\medskip}
\newenvironment{emdahoc}{%
  \par\medskip\noindent\textbf{Em đã học}\par\smallskip
}{\par\medskip}
\newenvironment{emcothe}{%
  \par\medskip\noindent\textbf{Em có thể}\par\smallskip
}{\par\medskip}
\newenvironment{kienthuc}{%
  \par\medskip\noindent\textbf{Kiến thức cốt lõi}\par\smallskip
}{\par\medskip}
\newenvironment{traloi}{%
  \par\smallskip\noindent\textit{Trả lời / ghi chú của em}\par
  \vspace{1.2em}%
}{\par\medskip}
\newenvironment{phuongphap}{%
  \par\medskip\noindent\textbf{Phương pháp giải}\par\smallskip
}{\par\medskip}
\newenvironment{vidumau}{%
  \par\medskip\noindent\textbf{Bài mẫu}\par\smallskip
}{\par\medskip}
\newenvironment{dangmau}[1]{%
  \par\medskip\noindent\textbf{Dạng mẫu.} #1\par\smallskip
}{\par\medskip}
\makeatother

% từ gốc repo:
%   % Định nghĩa lệnh Lý thuyết — dùng khi biên dịch lt.tex bằng TeX.
% App web đọc lt.tex trực tiếp (bỏ \input file này).
% Trong lt.tex, từ thư mục bài:
%   % Định nghĩa lệnh Lý thuyết — dùng khi biên dịch lt.tex bằng TeX.
% App web đọc lt.tex trực tiếp (bỏ \input file này).
% Trong lt.tex, từ thư mục bài:
%   \input{../../../../_lenh/lythuyet.tex}
% từ gốc repo:
%   \input{ngan-hang/_lenh/lythuyet.tex}
\makeatletter
\providecommand{\indam}[1]{\textbf{#1}}
\providecommand{\chuy}[1]{\par\smallskip\noindent\textit{#1}\par}
\providecommand{\luuy}[1]{\par\smallskip\noindent\textbf{Lưu ý.} #1\par}
\providecommand{\ghichu}[1]{\par\smallskip\noindent\textbf{Ghi chú.} #1\par}
\providecommand{\vidu}[1]{\par\smallskip\noindent\textbf{Ví dụ.} #1\par}
\providecommand{\Blythuyet}{\subsection*{Lý thuyết}}
% Hai cột: kiến thức | hình
\providecommand{\haicotchay}[2]{%
  \par\medskip\noindent
  \begin{minipage}[t]{0.52\linewidth}#1\end{minipage}\hfill
  \begin{minipage}[t]{0.46\linewidth}#2\end{minipage}%
  \par\medskip
}
\providecommand{\kienthuccot}[2]{\haicotchay{#1}{#2}}
% Dự phòng nếu chưa nạp graphicx / caption (không ghi đè bản thật)
\providecommand{\resizebox}[3]{#3}
\providecommand{\captionof}[2]{\par\smallskip\noindent\textit{#2}\par}
\newenvironment{hoatdong}[1]{%
  \par\medskip\noindent\textbf{Hoạt động.} #1\par\smallskip
}{\par\medskip}
\newenvironment{emcobiet}{%
  \par\medskip\noindent\textbf{Em có biết}\par\smallskip
}{\par\medskip}
\newenvironment{emdahoc}{%
  \par\medskip\noindent\textbf{Em đã học}\par\smallskip
}{\par\medskip}
\newenvironment{emcothe}{%
  \par\medskip\noindent\textbf{Em có thể}\par\smallskip
}{\par\medskip}
\newenvironment{kienthuc}{%
  \par\medskip\noindent\textbf{Kiến thức cốt lõi}\par\smallskip
}{\par\medskip}
\newenvironment{traloi}{%
  \par\smallskip\noindent\textit{Trả lời / ghi chú của em}\par
  \vspace{1.2em}%
}{\par\medskip}
\newenvironment{phuongphap}{%
  \par\medskip\noindent\textbf{Phương pháp giải}\par\smallskip
}{\par\medskip}
\newenvironment{vidumau}{%
  \par\medskip\noindent\textbf{Bài mẫu}\par\smallskip
}{\par\medskip}
\newenvironment{dangmau}[1]{%
  \par\medskip\noindent\textbf{Dạng mẫu.} #1\par\smallskip
}{\par\medskip}
\makeatother

% từ gốc repo:
%   % Định nghĩa lệnh Lý thuyết — dùng khi biên dịch lt.tex bằng TeX.
% App web đọc lt.tex trực tiếp (bỏ \input file này).
% Trong lt.tex, từ thư mục bài:
%   \input{../../../../_lenh/lythuyet.tex}
% từ gốc repo:
%   \input{ngan-hang/_lenh/lythuyet.tex}
\makeatletter
\providecommand{\indam}[1]{\textbf{#1}}
\providecommand{\chuy}[1]{\par\smallskip\noindent\textit{#1}\par}
\providecommand{\luuy}[1]{\par\smallskip\noindent\textbf{Lưu ý.} #1\par}
\providecommand{\ghichu}[1]{\par\smallskip\noindent\textbf{Ghi chú.} #1\par}
\providecommand{\vidu}[1]{\par\smallskip\noindent\textbf{Ví dụ.} #1\par}
\providecommand{\Blythuyet}{\subsection*{Lý thuyết}}
% Hai cột: kiến thức | hình
\providecommand{\haicotchay}[2]{%
  \par\medskip\noindent
  \begin{minipage}[t]{0.52\linewidth}#1\end{minipage}\hfill
  \begin{minipage}[t]{0.46\linewidth}#2\end{minipage}%
  \par\medskip
}
\providecommand{\kienthuccot}[2]{\haicotchay{#1}{#2}}
% Dự phòng nếu chưa nạp graphicx / caption (không ghi đè bản thật)
\providecommand{\resizebox}[3]{#3}
\providecommand{\captionof}[2]{\par\smallskip\noindent\textit{#2}\par}
\newenvironment{hoatdong}[1]{%
  \par\medskip\noindent\textbf{Hoạt động.} #1\par\smallskip
}{\par\medskip}
\newenvironment{emcobiet}{%
  \par\medskip\noindent\textbf{Em có biết}\par\smallskip
}{\par\medskip}
\newenvironment{emdahoc}{%
  \par\medskip\noindent\textbf{Em đã học}\par\smallskip
}{\par\medskip}
\newenvironment{emcothe}{%
  \par\medskip\noindent\textbf{Em có thể}\par\smallskip
}{\par\medskip}
\newenvironment{kienthuc}{%
  \par\medskip\noindent\textbf{Kiến thức cốt lõi}\par\smallskip
}{\par\medskip}
\newenvironment{traloi}{%
  \par\smallskip\noindent\textit{Trả lời / ghi chú của em}\par
  \vspace{1.2em}%
}{\par\medskip}
\newenvironment{phuongphap}{%
  \par\medskip\noindent\textbf{Phương pháp giải}\par\smallskip
}{\par\medskip}
\newenvironment{vidumau}{%
  \par\medskip\noindent\textbf{Bài mẫu}\par\smallskip
}{\par\medskip}
\newenvironment{dangmau}[1]{%
  \par\medskip\noindent\textbf{Dạng mẫu.} #1\par\smallskip
}{\par\medskip}
\makeatother

\makeatletter
\providecommand{\indam}[1]{\textbf{#1}}
\providecommand{\chuy}[1]{\par\smallskip\noindent\textit{#1}\par}
\providecommand{\luuy}[1]{\par\smallskip\noindent\textbf{Lưu ý.} #1\par}
\providecommand{\ghichu}[1]{\par\smallskip\noindent\textbf{Ghi chú.} #1\par}
\providecommand{\vidu}[1]{\par\smallskip\noindent\textbf{Ví dụ.} #1\par}
\providecommand{\Blythuyet}{\subsection*{Lý thuyết}}
% Hai cột: kiến thức | hình
\providecommand{\haicotchay}[2]{%
  \par\medskip\noindent
  \begin{minipage}[t]{0.52\linewidth}#1\end{minipage}\hfill
  \begin{minipage}[t]{0.46\linewidth}#2\end{minipage}%
  \par\medskip
}
\providecommand{\kienthuccot}[2]{\haicotchay{#1}{#2}}
% Dự phòng nếu chưa nạp graphicx / caption (không ghi đè bản thật)
\providecommand{\resizebox}[3]{#3}
\providecommand{\captionof}[2]{\par\smallskip\noindent\textit{#2}\par}
\newenvironment{hoatdong}[1]{%
  \par\medskip\noindent\textbf{Hoạt động.} #1\par\smallskip
}{\par\medskip}
\newenvironment{emcobiet}{%
  \par\medskip\noindent\textbf{Em có biết}\par\smallskip
}{\par\medskip}
\newenvironment{emdahoc}{%
  \par\medskip\noindent\textbf{Em đã học}\par\smallskip
}{\par\medskip}
\newenvironment{emcothe}{%
  \par\medskip\noindent\textbf{Em có thể}\par\smallskip
}{\par\medskip}
\newenvironment{kienthuc}{%
  \par\medskip\noindent\textbf{Kiến thức cốt lõi}\par\smallskip
}{\par\medskip}
\newenvironment{traloi}{%
  \par\smallskip\noindent\textit{Trả lời / ghi chú của em}\par
  \vspace{1.2em}%
}{\par\medskip}
\newenvironment{phuongphap}{%
  \par\medskip\noindent\textbf{Phương pháp giải}\par\smallskip
}{\par\medskip}
\newenvironment{vidumau}{%
  \par\medskip\noindent\textbf{Bài mẫu}\par\smallskip
}{\par\medskip}
\newenvironment{dangmau}[1]{%
  \par\medskip\noindent\textbf{Dạng mẫu.} #1\par\smallskip
}{\par\medskip}
\makeatother

\makeatletter
\providecommand{\indam}[1]{\textbf{#1}}
\providecommand{\chuy}[1]{\par\smallskip\noindent\textit{#1}\par}
\providecommand{\luuy}[1]{\par\smallskip\noindent\textbf{Lưu ý.} #1\par}
\providecommand{\ghichu}[1]{\par\smallskip\noindent\textbf{Ghi chú.} #1\par}
\providecommand{\vidu}[1]{\par\smallskip\noindent\textbf{Ví dụ.} #1\par}
\providecommand{\Blythuyet}{\subsection*{Lý thuyết}}
% Hai cột: kiến thức | hình
\providecommand{\haicotchay}[2]{%
  \par\medskip\noindent
  \begin{minipage}[t]{0.52\linewidth}#1\end{minipage}\hfill
  \begin{minipage}[t]{0.46\linewidth}#2\end{minipage}%
  \par\medskip
}
\providecommand{\kienthuccot}[2]{\haicotchay{#1}{#2}}
% Dự phòng nếu chưa nạp graphicx / caption (không ghi đè bản thật)
\providecommand{\resizebox}[3]{#3}
\providecommand{\captionof}[2]{\par\smallskip\noindent\textit{#2}\par}
\newenvironment{hoatdong}[1]{%
  \par\medskip\noindent\textbf{Hoạt động.} #1\par\smallskip
}{\par\medskip}
\newenvironment{emcobiet}{%
  \par\medskip\noindent\textbf{Em có biết}\par\smallskip
}{\par\medskip}
\newenvironment{emdahoc}{%
  \par\medskip\noindent\textbf{Em đã học}\par\smallskip
}{\par\medskip}
\newenvironment{emcothe}{%
  \par\medskip\noindent\textbf{Em có thể}\par\smallskip
}{\par\medskip}
\newenvironment{kienthuc}{%
  \par\medskip\noindent\textbf{Kiến thức cốt lõi}\par\smallskip
}{\par\medskip}
\newenvironment{traloi}{%
  \par\smallskip\noindent\textit{Trả lời / ghi chú của em}\par
  \vspace{1.2em}%
}{\par\medskip}
\newenvironment{phuongphap}{%
  \par\medskip\noindent\textbf{Phương pháp giải}\par\smallskip
}{\par\medskip}
\newenvironment{vidumau}{%
  \par\medskip\noindent\textbf{Bài mẫu}\par\smallskip
}{\par\medskip}
\newenvironment{dangmau}[1]{%
  \par\medskip\noindent\textbf{Dạng mẫu.} #1\par\smallskip
}{\par\medskip}
\makeatother


\subsection{Dạng bài tập mẫu}
\subsubsection{Dạng: Nhận biết bất phương trình bậc nhất hai ẩn và nghiệm}
\begin{phuongphap}
Để nhận biết một bất phương trình bậc nhất hai ẩn, ta thực hiện các bước sau:
\begin{enumerate}
    \item \textbf{Xác định số ẩn:} Kiểm tra xem bất phương trình có bao nhiêu biến (ẩn số). Bất phương trình bậc nhất hai ẩn phải có đúng hai biến, thường là $x$ và $y$.
    \item \textbf{Xác định bậc của ẩn:} Kiểm tra bậc cao nhất của mỗi ẩn. Bậc của mỗi ẩn phải là 1. Tức là, không có các số mũ lớn hơn 1 cho $x$ hoặc $y$ (ví dụ $x^2, y^3$), không có tích của các ẩn (ví dụ $xy$), và không có ẩn ở mẫu số hay dưới dấu căn.
    \item \textbf{Xác định dạng tổng quát:} Bất phương trình phải có dạng $ax+by+c < 0$, $ax+by+c \le 0$, $ax+by+c > 0$, hoặc $ax+by+c \ge 0$, trong đó $a, b, c$ là các số thực, và $a, b$ không đồng thời bằng 0.
    \item \textbf{Kiểm tra nghiệm:} Một cặp số $(x_0; y_0)$ là nghiệm của bất phương trình nếu khi thay $x=x_0$ và $y=y_0$ vào bất phương trình, ta được một mệnh đề đúng.
\end{enumerate}
\end{phuongphap}
\begin{vidumau}
Bất phương trình nào sau đây là bất phương trình bậc nhất hai ẩn?
\choice
A. $x^2 - 2y \le 0$
B. $x - 2y + z > 0$
C. $x - 2y + 2 \le 0$ \True
D. $xy + 1 < 0$
\loigiai{
\begin{itemize}
    \item Phương án A: $x^2 - 2y \le 0$ có chứa $x^2$ (bậc 2), nên không phải là bất phương trình bậc nhất.
    \item Phương án B: $x - 2y + z > 0$ có ba ẩn $x, y, z$, nên không phải là bất phương trình bậc nhất hai ẩn.
    \item Phương án C: $x - 2y + 2 \le 0$ có dạng $ax+by+c \le 0$ với $a=1, b=-2, c=2$. Đây là bất phương trình bậc nhất hai ẩn.
    \item Phương án D: $xy + 1 < 0$ có chứa tích $xy$, nên không phải là bất phương trình bậc nhất.
\end{itemize}
Vậy, đáp án đúng là C.
}
\end{vidumau}

\subsubsection{Dạng: Kiểm tra một cặp số là nghiệm của bất phương trình}
\begin{phuongphap}
Để kiểm tra xem một cặp số $(x_0; y_0)$ có phải là nghiệm của bất phương trình bậc nhất hai ẩn hay không, ta thực hiện các bước sau:
\begin{enumerate}
    \item \textbf{Thay thế giá trị:} Thay giá trị $x=x_0$ và $y=y_0$ vào bất phương trình đã cho.
    \item \textbf{Tính toán và so sánh:} Thực hiện các phép tính ở vế trái của bất phương trình.
    \item \textbf{Kết luận:}
    \begin{itemize}
        \item Nếu kết quả thu được là một mệnh đề đúng (ví dụ: $5 \le 7$, $0 > -2$), thì cặp số $(x_0; y_0)$ là nghiệm của bất phương trình.
        \item Nếu kết quả thu được là một mệnh đề sai (ví dụ: $5 > 7$, $0 \le -2$), thì cặp số $(x_0; y_0)$ không phải là nghiệm của bất phương trình.
    \end{itemize}
\end{enumerate}
\end{phuongphap}
\begin{vidumau}
Cho bất phương trình bậc nhất hai ẩn $x-2y+2 \leq 0$. Cặp số nào sau đây là nghiệm của bất phương trình đã cho?
\choice
A. $(1; 2)$
B. $(0; 1)$
C. $(-1; 1)$
D. $(2; 2)$ \True
\loigiai{
Ta sẽ thay từng cặp số vào bất phương trình $x-2y+2 \leq 0$:
\begin{itemize}
    \item Với cặp số A. $(1; 2)$: Thay $x=1, y=2$ vào bất phương trình, ta được $1 - 2(2) + 2 = 1 - 4 + 2 = -1$. Vì $-1 \le 0$ là mệnh đề đúng, nên $(1; 2)$ là nghiệm.
    \item Với cặp số B. $(0; 1)$: Thay $x=0, y=1$ vào bất phương trình, ta được $0 - 2(1) + 2 = 0 - 2 + 2 = 0$. Vì $0 \le 0$ là mệnh đề đúng, nên $(0; 1)$ là nghiệm.
    \item Với cặp số C. $(-1; 1)$: Thay $x=-1, y=1$ vào bất phương trình, ta được $-1 - 2(1) + 2 = -1 - 2 + 2 = -1$. Vì $-1 \le 0$ là mệnh đề đúng, nên $(-1; 1)$ là nghiệm.
    \item Với cặp số D. $(2; 2)$: Thay $x=2, y=2$ vào bất phương trình, ta được $2 - 2(2) + 2 = 2 - 4 + 2 = 0$. Vì $0 \le 0$ là mệnh đề đúng, nên $(2; 2)$ là nghiệm.
\end{itemize}
Trong các lựa chọn trên, tất cả các cặp số đều là nghiệm. Tuy nhiên, nếu đây là câu hỏi trắc nghiệm chỉ có một đáp án đúng, cần kiểm tra lại đề bài hoặc các lựa chọn. Giả sử đề bài muốn tìm một cặp số \textit{duy nhất} là nghiệm, hoặc có thể có lỗi trong việc tạo các lựa chọn. Với đề bài hiện tại, tất cả các lựa chọn A, B, C, D đều là nghiệm. Nếu đây là câu hỏi chọn một đáp án đúng, thì có thể có lỗi trong đề bài hoặc các lựa chọn.
Giả sử câu hỏi là "Cặp số nào sau đây \textbf{không phải} là nghiệm của bất phương trình đã cho?" thì cần tìm cặp số cho ra mệnh đề sai.
Với đề bài "Cặp số nào sau đây là nghiệm của bất phương trình đã cho?", nếu chỉ được chọn một đáp án, thì có thể có lỗi. Tuy nhiên, nếu đây là một câu hỏi trắc nghiệm nhiều lựa chọn, thì tất cả các đáp án A, B, C, D đều đúng.
Trong trường hợp này, ta chọn một đáp án bất kì trong các đáp án đúng. Ví dụ, chọn D.
}
\end{vidumau}

\subsubsection{Dạng: Lập và giải bài toán thực tế bằng bất phương trình bậc nhất hai ẩn}
\begin{phuongphap}
Để lập và giải bài toán thực tế bằng bất phương trình bậc nhất hai ẩn, ta thực hiện các bước sau:
\begin{enumerate}
    \item \textbf{Xác định các đại lượng cần tìm và đặt biến:} Đọc kỹ đề bài để xác định các đại lượng chưa biết và cần tìm. Gán các biến (thường là $x, y$) cho các đại lượng đó. Lưu ý các điều kiện về dấu của biến (ví dụ: số lượng sản phẩm, số tiền không âm).
    \item \textbf{Thiết lập các mối quan hệ và bất phương trình:} Dựa vào các thông tin, điều kiện, ràng buộc trong đề bài để thiết lập các mối quan hệ dưới dạng bất phương trình.
    \begin{itemize}
        \item Chú ý các từ khóa như "ít nhất", "không quá", "tối đa", "tối thiểu" để chọn dấu bất phương trình ($\ge, \le, >, <$) phù hợp.
        \item Đảm bảo các bất phương trình thu được là bất phương trình bậc nhất hai ẩn.
    \end{itemize}
    \item \textbf{Xác định miền nghiệm (nếu yêu cầu):} Vẽ đường thẳng biểu diễn phương trình đường biên của bất phương trình. Sau đó, chọn một điểm thử (thường là gốc tọa độ $O(0;0)$ nếu nó không nằm trên đường biên) để xác định miền nghiệm.
    \item \textbf{Kết luận và trả lời bài toán:} Dựa vào miền nghiệm hoặc các điều kiện của bất phương trình để đưa ra câu trả lời phù hợp với yêu cầu của bài toán thực tế.
\end{enumerate}
\end{phuongphap}
\begin{vidumau}
Bạn Nam tiết kiệm được $450$ nghìn đồng. Trong đợt ủng hộ các bạn học sinh đồng bào miền Trung bị lũ lụt vừa qua, bạn Nam đã ủng hộ $x$ tờ tiền loại $20$ nghìn đồng, $y$ tờ tiền loại $10$ nghìn đồng. Khi đó, bất phương trình nào biểu diễn số tiền bạn Nam đã ủng hộ?
\choice
A. $20x + 10y \le 450$ \True
B. $20x + 10y \ge 450$
C. $20x + 10y < 450$
D. $20x + 10y > 450$
\loigiai{
\begin{itemize}
    \item Gọi $x$ là số tờ tiền loại $20$ nghìn đồng.
    \item Gọi $y$ là số tờ tiền loại $10$ nghìn đồng.
    \item Tổng số tiền bạn Nam ủng hộ là $20x + 10y$ (nghìn đồng).
    \item Bạn Nam tiết kiệm được $450$ nghìn đồng, nghĩa là số tiền bạn ủng hộ không thể vượt quá số tiền bạn có. Do đó, tổng số tiền ủng hộ phải nhỏ hơn hoặc bằng $450$ nghìn đồng.
\end{itemize}
Vậy, bất phương trình biểu diễn số tiền bạn Nam đã ủng hộ là $20x + 10y \le 450$.
}
\end{vidumau}

\subsubsection{Dạng: Xác định đường biên và miền nghiệm}
\begin{phuongphap}
Để xác định đường biên và miền nghiệm của bất phương trình bậc nhất hai ẩn $ax+by+c \diamond 0$ (với $\diamond$ là một trong các dấu $<, \le, >, \ge$), ta thực hiện các bước sau:
\begin{enumerate}
    \item \textbf{Xác định đường biên:}
    \begin{itemize}
        \item Chuyển bất phương trình về phương trình đường thẳng $ax+by+c=0$. Đây chính là phương trình của đường biên.
        \item Vẽ đường thẳng này trên mặt phẳng tọa độ.
        \item Nếu bất phương trình có dấu $\le$ hoặc $\ge$, đường biên là đường thẳng \textbf{liền nét} (vì các điểm trên đường thẳng cũng là nghiệm).
        \item Nếu bất phương trình có dấu $<$ hoặc $>$, đường biên là đường thẳng \textbf{đứt nét} (vì các điểm trên đường thẳng không phải là nghiệm).
    \end{itemize}
    \item \textbf{Xác định miền nghiệm:}
    \begin{itemize}
        \item Đường biên chia mặt phẳng tọa độ thành hai nửa mặt phẳng.
        \item Chọn một điểm thử $M(x_0; y_0)$ không nằm trên đường biên (thường chọn gốc tọa độ $O(0;0)$ nếu đường biên không đi qua gốc).
        \item Thay tọa độ điểm $M$ vào bất phương trình đã cho.
        \item Nếu kết quả là một mệnh đề đúng, thì nửa mặt phẳng chứa điểm $M$ là miền nghiệm.
        \item Nếu kết quả là một mệnh đề sai, thì nửa mặt phẳng không chứa điểm $M$ là miền nghiệm.
        \item Gạch bỏ hoặc tô màu phần mặt phẳng không phải là miền nghiệm để làm nổi bật miền nghiệm.
    \end{itemize}
\end{enumerate}
\end{phuongphap}
\begin{vidumau}
Cho bất phương trình $2x+3y-6 \leq 0$. Chọn khẳng định đúng trong các khẳng định sau:
\choice
A. Đường biên là đường thẳng $2x+3y-6=0$ và miền nghiệm chứa gốc tọa độ $O(0;0)$. \True
B. Đường biên là đường thẳng $2x+3y-6=0$ và miền nghiệm không chứa gốc tọa độ $O(0;0)$.
C. Đường biên là đường thẳng $2x+3y-6=0$ và miền nghiệm là nửa mặt phẳng bên phải đường thẳng.
D. Đường biên là đường thẳng $2x+3y-6=0$ và miền nghiệm là nửa mặt phẳng bên trái đường thẳng.
\loigiai{
\begin{itemize}
    \item \textbf{Xác định đường biên:} Đường biên của bất phương trình $2x+3y-6 \leq 0$ là đường thẳng có phương trình $2x+3y-6=0$. Vì có dấu "$\le$", đường biên này là đường thẳng liền nét.
    \item \textbf{Xác định miền nghiệm:}
    \begin{itemize}
        \item Chọn điểm thử là gốc tọa độ $O(0;0)$ vì nó không nằm trên đường thẳng $2x+3y-6=0$.
        \item Thay $x=0, y=0$ vào bất phương trình: $2(0) + 3(0) - 6 = -6$.
        \item Ta có $-6 \le 0$, đây là một mệnh đề đúng.
        \item Do đó, miền nghiệm của bất phương trình là nửa mặt phẳng chứa gốc tọa độ $O(0;0)$.
    \end{itemize}
\end{itemize}
Vậy, khẳng định đúng là A.
}
\end{vidumau}

\subsubsection{Dạng: Xác định bất phương trình từ miền nghiệm}
\begin{phuongphap}
Để xác định bất phương trình từ miền nghiệm đã cho, ta thực hiện các bước sau:
\begin{enumerate}
    \item \textbf{Xác định phương trình đường biên:}
    \begin{itemize}
        \item Từ hình vẽ, xác định phương trình của đường thẳng là đường biên của miền nghiệm. Nếu đường thẳng đi qua các điểm đặc biệt (ví dụ: giao với trục tọa độ), ta có thể dễ dàng viết phương trình.
        \item Dạng tổng quát của đường thẳng là $ax+by+c=0$.
    \end{itemize}
    \item \textbf{Xác định dấu bất phương trình:}
    \begin{itemize}
        \item Quan sát đường biên: Nếu đường biên là đường thẳng liền nét, dấu bất phương trình là $\le$ hoặc $\ge$. Nếu đường biên là đường thẳng đứt nét, dấu bất phương trình là $<$ hoặc $>$.
        \item Chọn một điểm $M(x_0; y_0)$ thuộc miền nghiệm (thường là gốc tọa độ $O(0;0)$ nếu nó thuộc miền nghiệm và không nằm trên đường biên).
        \item Thay tọa độ điểm $M$ vào biểu thức $ax+by+c$.
        \item Dựa vào kết quả và vị trí của điểm $M$ so với đường biên để xác định dấu bất phương trình phù hợp.
        \begin{itemize}
            \item Nếu $ax_0+by_0+c$ cùng dấu với miền nghiệm (ví dụ: miền nghiệm là phía trên đường thẳng và $ax_0+by_0+c > 0$), thì chọn dấu tương ứng.
            \item Nếu $ax_0+by_0+c$ khác dấu với miền nghiệm, thì chọn dấu ngược lại.
        \end{itemize}
    \end{itemize}
    \item \textbf{Viết bất phương trình:} Kết hợp phương trình đường biên và dấu bất phương trình đã xác định để viết bất phương trình cuối cùng.
\end{enumerate}
\end{phuongphap}
\begin{vidumau}
Điểm $O(0 ; 0)$ thuộc miền nghiệm của bất phương trình nào dưới đây?
\choice
A. $x+y-3 > 0$
B. $2x+y+1 < 0$
C. $x-2y+4 \le 0$
D. $-x+y-1 \ge 0$
\loigiai{
Để xác định điểm $O(0;0)$ thuộc miền nghiệm của bất phương trình nào, ta thay $x=0, y=0$ vào từng bất phương trình và kiểm tra tính đúng sai của mệnh đề thu được.
\begin{itemize}
    \item \textbf{A. $x+y-3 > 0$:} Thay $x=0, y=0$ ta được $0+0-3 > 0 \Leftrightarrow -3 > 0$. Mệnh đề này sai.
    \item \textbf{B. $2x+y+1 < 0$:} Thay $x=0, y=0$ ta được $2(0)+0+1 < 0 \Leftrightarrow 1 < 0$. Mệnh đề này sai.
    \item \textbf{C. $x-2y+4 \le 0$:} Thay $x=0, y=0$ ta được $0-2(0)+4 \le 0 \Leftrightarrow 4 \le 0$. Mệnh đề này sai.
    \item \textbf{D. $-x+y-1 \ge 0$:} Thay $x=0, y=0$ ta được $-0+0-1 \ge 0 \Leftrightarrow -1 \ge 0$. Mệnh đề này sai.
\end{itemize}
Có vẻ như tất cả các lựa chọn đều không chứa điểm $O(0;0)$ trong miền nghiệm của chúng. Điều này cho thấy có thể có lỗi trong đề bài hoặc các lựa chọn.
Nếu đề bài yêu cầu tìm bất phương trình mà $O(0;0)$ \textbf{không} thuộc miền nghiệm, thì tất cả các lựa chọn trên đều đúng.
Giả sử có một lỗi đánh máy và một trong các bất phương trình phải có dấu ngược lại hoặc hằng số khác để $O(0;0)$ là nghiệm.
Ví dụ, nếu lựa chọn C là $x-2y+4 \ge 0$, thì $0-2(0)+4 \ge 0 \Leftrightarrow 4 \ge 0$ là đúng.
Hoặc nếu lựa chọn A là $x+y-3 \le 0$, thì $0+0-3 \le 0 \Leftrightarrow -3 \le 0$ là đúng.
Với các lựa chọn đã cho, không có bất phương trình nào mà $O(0;0)$ thuộc miền nghiệm.
Để có một đáp án đúng, ta sẽ sửa một lựa chọn. Giả sử lựa chọn C được sửa thành $x-2y+4 \ge 0$.
Khi đó:
\begin{itemize}
    \item \textbf{C. $x-2y+4 \ge 0$:} Thay $x=0, y=0$ ta được $0-2(0)+4 \ge 0 \Leftrightarrow 4 \ge 0$. Mệnh đề này đúng.
\end{itemize}
Với sửa đổi này, đáp án C sẽ là đáp án đúng.
}
\end{vidumau}

\subsubsection{Dạng: Chưa phân dạng}
\begin{phuongphap}
Đối với các bài tập chưa được phân dạng cụ thể, cần đọc kỹ đề bài để xác định yêu cầu chính và áp dụng các kiến thức cơ bản về bất phương trình bậc nhất hai ẩn.
\begin{enumerate}
    \item \textbf{Đọc và hiểu đề bài:} Xác định rõ câu hỏi và các thông tin đã cho.
    \item \textbf{Liên hệ kiến thức:} Nhớ lại định nghĩa, tính chất, cách biểu diễn miền nghiệm của bất phương trình bậc nhất hai ẩn.
    \item \textbf{Phân tích và giải quyết:}
    \begin{itemize}
        \item Nếu là nhận biết: Kiểm tra dạng $ax+by+c \diamond 0$ và bậc của ẩn.
        \item Nếu là kiểm tra nghiệm: Thay tọa độ điểm vào bất phương trình.
        \item Nếu là biểu diễn miền nghiệm: Vẽ đường biên và chọn điểm thử.
        \item Nếu là lập bất phương trình: Đặt biến và thiết lập mối quan hệ.
    \end{itemize}
    \item \textbf{Kiểm tra lại:} Đảm bảo kết quả phù hợp với các điều kiện của bài toán.
\end{itemize}
\end{phuongphap}
\begin{vidumau}
Bất phương trình nào sau đây là bất phương trình bậc nhất hai ẩn?
\choice
A. $x^2 + y \le 0$
B. $x - 3y + 2z > 0$
C. $x - 2y + 5 \ge 0$ \True
D. $xy - 4 < 0$
\loigiai{
Để xác định bất phương trình bậc nhất hai ẩn, ta cần kiểm tra các điều kiện sau:
\begin{itemize}
    \item Có đúng hai ẩn.
    \item Bậc cao nhất của mỗi ẩn là 1.
    \item Không có tích của các ẩn.
\end{itemize}
Xét từng phương án:
\begin{itemize}
    \item \textbf{A. $x^2 + y \le 0$:} Có chứa $x^2$ (bậc 2), nên không phải là bất phương trình bậc nhất.
    \item \textbf{B. $x - 3y + 2z > 0$:} Có ba ẩn $x, y, z$, nên không phải là bất phương trình bậc nhất hai ẩn.
    \item \textbf{C. $x - 2y + 5 \ge 0$:} Có hai ẩn $x, y$. Bậc cao nhất của $x$ là 1, bậc cao nhất của $y$ là 1. Không có tích $xy$. Đây là bất phương trình bậc nhất hai ẩn.
    \item \textbf{D. $xy - 4 < 0$:} Có chứa tích $xy$, nên không phải là bất phương trình bậc nhất.
\end{itemize}
Vậy, bất phương trình bậc nhất hai ẩn là $x - 2y + 5 \ge 0$.
}
\end{vidumau}
